% =====================================================================
%  SUPPLEMENTARY MATERIAL
%  "Contagion: When Per-Hop Prompt-Injection Survival in LLM Agent Networks
%   Does (and Does Not) Compose"
%
%  Chứa Appendices A-D được tham chiếu từ bài chính.
%  Biên dịch:  pdflatex supplementary.tex   (chạy 2 lần)
%
%  LƯU Ý: file này KHÔNG dùng \ref tới file chính (không phụ thuộc .aux của bài
%  chính). Mọi tham chiếu tới bài chính đều ghi thẳng số mục, ví dụ
%  "Section 3.8 of the main paper".
% =====================================================================
\documentclass[11pt,a4paper]{article}

\usepackage{amsmath,amssymb}
\usepackage{booktabs}
\usepackage{graphicx}
\usepackage[margin=2.4cm]{geometry}
\usepackage[hidelinks]{hyperref}

% \ind là macro của bài chính (indicator function); phụ lục không \input bài chính
% nên phải định nghĩa lại, nếu không sẽ lỗi "Undefined control sequence".
\newcommand{\ind}{\mathbf{1}}

\title{\vspace{-1.5cm}\textbf{Supplementary Material}\\[0.4em]
\large Contagion: When Per-Hop Prompt-Injection Survival in\\ LLM Agent Networks Does (and Does Not) Compose}
\author{Anonymous Author(s)}
\date{}

\begin{document}
\maketitle
\vspace{-1.2cm}



\appendix

% =====================================================================
\section{Multiplicity Control for the Composition Test}
\label{app:bh}

Table~\ref{tab:bh} applies a Benjamini--Hochberg correction to the composition
tests of Section~3.7 of the main paper, using the \emph{default fresh-artefact
protocol} for every cell at the common $n = 40$ grid. Only the Llama cell rejects
$H_0$; it survives the correction at $q = 0.05$, while the four remaining
non-degenerate cells---including DeepSeek, whose fixed-artefact deviation is a
protocol artefact that disappears under the fresh protocol (Section~5.3 of the main
paper)---stay non-significant. The Claude cell is degenerate
($\mathrm{ASR} = \prod_i \hat s_i = 0$, $\mathrm{sd}(\Delta) = 0$) and is excluded
from the correction. At the higher $n = 200$ resolution the Nova cell also crosses
into rejection (Section~5.10 of the main paper); we keep the correction at the
common $n = 40$ grid so the five cells are compared at equal power.

\begin{table}[h]
\small
\centering
\caption{Benjamini--Hochberg over the $m = 5$ non-degenerate chain cells
(fresh-artefact protocol). $p$-values are the bootstrap two-sided values from
Section~3.7 of the main paper. Only the Llama rejection survives.}
\vspace{10pt} % <--- Đã sửa thành 10pt
\label{tab:bh}
\begin{tabular}{clcc}
\toprule
\textit{rank} & \textit{cell} & \textit{$p$} & \textit{BH thr.\ ($q{=}0.05$)} \\
\midrule
1 & Llama $\cdot$ none            & $0.0002$ & $0.010$ \\
2 & Nova $\cdot$ none             & $0.0840$ & $0.020$ \\
3 & qwen $\cdot$ none             & $0.3165$ & $0.030$ \\
4 & DeepSeek $\cdot$ none (fresh) & $0.6300$ & $0.040$ \\
5 & qwen $\cdot$ paraphrase       & $0.7775$ & $0.050$ \\
\bottomrule
\end{tabular}
\end{table}

% =====================================================================
\section{Context-Clustered Interval on the Weak Edge}
\label{app:cluster}

Section~5.10 of the main paper reports that a homogeneity test rejects exchangeability
across benign task contexts on the weak Llama edge. Because that violates the
i.i.d.\ assumption behind Wilson intervals, we recompute the edge rate with a
cluster bootstrap that resamples \emph{repeats} (seed replications), not individual
trials. At a fixed temperature ($T = 0.7$, six repeats of twenty trials) the two
intervals nearly coincide (Table~\ref{tab:cluster}): the design-effect widening is
$0.99\times$, so within a temperature the Wilson interval is well calibrated. The
overdispersion $\varphi = 2.93$ reported in Section~4 of the main paper is a
\emph{cross-temperature} artefact---temperature is a deliberate systematic factor,
not trial noise---which is why we compute $\varphi$ within a temperature.

\begin{table}[h]
\small
\centering
\caption{Weak edge $\texttt{a1}\to\texttt{a2}$ (Llama, $T = 0.7$, $n = 120$):
i.i.d.\ Wilson versus a repeat-clustered bootstrap ($10^4$ resamples).}
\vspace{10pt} % <--- Đã sửa thành 10pt
\label{tab:cluster}
\begin{tabular}{lccc}
\toprule
\textit{interval} & \textit{$\hat s$} & \textit{95\% CI} & \textit{width} \\
\midrule
Wilson (i.i.d.)              & $0.242$ & $[0.174, 0.326]$ & $0.152$ \\
cluster bootstrap (by repeat) & $0.242$ & $[0.167, 0.317]$ & $0.150$ \\
\bottomrule
\end{tabular}
\end{table}

% =====================================================================
\section{Threshold Sensitivity of the Floor Effect}
\label{app:tau}

A referee may ask whether the ``resistant'' verdict on the frontier model is an
artefact of the compromise threshold $\tau_{\mathrm{ASV}} = 0.9$. We re-score the
single-hop obfuscation cells (Claude, $n = 30$ per cell) at
$\tau \in \{0.7, 0.8, 0.9, 1.0\}$ directly from the stored ASV scores---a
deterministic re-thresholding, no new model calls. The compromise rate is
\emph{invariant} across $\tau$ in every cell (Table~\ref{tab:tau}): the no-defence
baseline stays at $0.17$ and redaction at $0.00$ regardless of the threshold. The
reason is that the ASV distribution is bimodal---of $180$ trials, $158$ score $0.0$
and $14$ score $1.0$; the only eight intermediate trials all score $\le 0.667$, i.e.\
below the lowest threshold tested. The floor effect is therefore a property of the
model, not of where the judge's cut is placed.

\begin{table}[h]
\small
\centering
\caption{Single-hop compromise rate versus $\tau_{\mathrm{ASV}}$ (Claude, $n = 30$
per cell), re-scored offline from stored ASV. Rates are threshold-invariant.}
\vspace{10pt} % <--- Đã sửa thành 10pt (Bảng 3 trong hình)
\label{tab:tau}
\begin{tabular}{llcccc}
\toprule
\textit{defence} & \textit{form} & \textit{$\tau{=}0.7$} & \textit{$0.8$}
& \textit{$0.9$} & \textit{$1.0$} \\
\midrule
none   & plain      & $0.17$ & $0.17$ & $0.17$ & $0.17$ \\
none   & spaced     & $0.17$ & $0.17$ & $0.17$ & $0.17$ \\
none   & split-word & $0.13$ & $0.13$ & $0.13$ & $0.13$ \\
redact & plain      & $0.00$ & $0.00$ & $0.00$ & $0.00$ \\
redact & spaced     & $0.00$ & $0.00$ & $0.00$ & $0.00$ \\
redact & split-word & $0.00$ & $0.00$ & $0.00$ & $0.00$ \\
\bottomrule
\end{tabular}
\end{table}

% =====================================================================
\section{Theory Assumptions: Independence, Correlation, and $R_0$}
\label{app:theory-assumptions}

This appendix clarifies the theoretical assumptions invoked in Section 3.8 of the
main paper.

\paragraph{Product-of-complements and correlation.} The percolation recurrence
\[
\Pr[v] = 1 - \prod_{u \to v}\!\bigl(1 - s_{u\to v}\,\Pr[u]\bigr)
\]
assumes \emph{edge activations} (and the corresponding parent-reach events relevant to the
product form) are conditionally independent given the graph structure. When multiple
parents $u_1,\ldots,u_k$ of $v$ share a common ancestor, the events $\{\Pr[u_i]\}$ may be
correlated, so $\Pr[\bigwedge_i (\text{not reached via }u_i)] \neq \prod_i (1 - s_{u_i\to v}\Pr[u_i])$
in general. For the specific topologies evaluated in this work (chains, trees, star,
and the small manager/worker loop), we do not rely on that product form across shared
ancestors; where we use it, we treat it as the marginal approximation under the stated
conditions or restrict claims to graph classes where the independence holds. For
arbitrary DAGs with shared ancestors, a path-based (inclusion--exclusion) formulation
or explicit correlation modelling is required.

\paragraph{$R_0$ definition and scope.} For strictly acyclic (feed-forward) graphs, the
standard linearized next-generation matrix $M$ with $M_{uv}=s_{u\to v}$ is strictly
triangular in a topological ordering, hence $\rho(M)=0$. This is a structural property of
acyclicity, not an empirical safety guarantee. We therefore report $R_0$ as a
\emph{descriptive, finite-horizon} reproduction statistic: specifically, we define
$R_0$ as the mean number of newly compromised neighbours produced per compromised agent
instance over the observed finite horizon (generation convention as stated in
Section 4.2 of the main paper), with its denominator clearly specified. We do not
interpret $\rho(M)<1$ as a threshold for safety on finite acyclic graphs; reachability
is given by the percolation (Eq.~3 in the main paper), not by $R_0$ alone.

\paragraph{Parallel paths and defence placement.} Hardening a single node is not
equally effective when there exist multiple $s$-$t$ paths: effectiveness depends on
whether that node lies on all paths (a cut). For unique-path graphs the effect is
path-specific; we do not generalise equal effectiveness to multi-path DAGs. Placement
claims are limited to the graph classes we evaluate.

\paragraph{Transportability and product composition.} The chain identity
$\mathrm{ASR}=\prod_i s_i^{\mathrm{nat}}$ is algebraic given the Markov/chain
structure. However, composing \emph{isolated} per-edge estimates $s^{\mathrm{ctrl}}$
requires that $s^{\mathrm{ctrl}}_i \approx s^{\mathrm{nat}}_i$ (transportability). We do
not assume this holds in general; we test it empirically and use the depth-error slope
as a diagnostic (Section 5.2 of the main paper).

% =====================================================================
\section{Recurrence Makes the Threshold Non-Vacuous}
\label{app:cyclic}

This appendix gives the full treatment of the recurrent (cyclic) case summarised in
Section~5.8 of the main paper. We test a pattern real frameworks use: a manager
broadcasting to workers that report back.

For a manager with $w$ reporting workers, the infection matrix has
$M[m\!\to\! w_j] = a_j$ and $M[w_j\!\to\! m] = c_j$, and its characteristic
polynomial is $\lambda^{w-1}(\lambda^2 - \sum_j a_j c_j)$. Hence
\begin{equation}
  \rho(M) = \sqrt{\textstyle\sum_j a_j\,c_j}
  \label{eq:cyclicrho}
\end{equation}
for every $w$ (checked numerically against the eigenvalues to machine precision on
$320$ random configurations). Two consequences follow. A feed-forward reading sets
every $c_j = 0$, so $\rho = 0$ however strong the edges are---the vacuity of
Section~3.8 of the main paper again, now with the feedback edges removed. And the
system is supercritical exactly when $\sum_j a_j c_j > 1$; for symmetric edges
$a_j = c_j = s$ this is $w\,s^2 > 1$. That is a \emph{design rule} rather than a
diagnostic: at $s = 0.5$ a manager needs five reporting workers to sustain an
endemic injection, at $s = 0.7$ three, at $s = 0.9$ two.

Applied to measured survivals, Eq.~\eqref{eq:cyclicrho} makes opposite
predictions, which is the test a criterion has to pass to be worth using
(Table~\ref{tab:cyclic}). On Llama the sum is $1.800$ and on DeepSeek $1.018$, so both
frontier loops should be supercritical; on qwen2.5:7b it is $0.578$, so the same loop
should decay.

\begin{table}[h]
\small
\centering
\caption{One criterion, three models, opposite predictions. $\sum_j a_jc_j$ uses
the measured per-hop survivals of the same three-agent recurrent pattern; the
feed-forward column is the same survivals with the two feedback edges removed.}
\vspace{10pt} % <--- Đã sửa thành 10pt
\label{tab:cyclic}
\begin{tabular}{lcccc}
\toprule
\textit{model} & \textit{$\sum_j a_jc_j$} & \textit{$\rho$ rec.}
& \textit{$\rho$ fwd.} & \textit{prediction} \\
\midrule
Llama 3.3 70B & $1.800$ & $\mathbf{1.342}$ & $0.000$ & supercritical \\
DeepSeek V3.2 & $1.018$ & $\mathbf{1.009}$ & $0.000$ & supercritical (marginal) \\
qwen2.5:7b    & $0.578$ & $0.760$ & $0.000$ & subcritical \\
\bottomrule
\end{tabular}
\end{table}

The measured dynamics follow. On Llama the recurrent manager is re-compromised by
its own workers' reports in $0.725$ $[0.572, 0.839]$ of runs and does not die out:
at round five it is still compromised in $0.700$ of runs with the workers near
saturation ($0.925$, $0.975$), as $\rho > 1$ predicts; extending the same loop to
eight rounds leaves it endemic (manager $0.583$, workers $0.900$). DeepSeek, whose
$\sum_j a_jc_j = 1.018$ sits just above criticality, behaves the same way over eight
rounds---the manager settles at $0.567$ with workers at $0.767$ and $0.700$ rather
than decaying---so the recurrent verdict is not a single-model artefact. On qwen the
same loop decays: the manager's re-compromise rate falls from $0.775$
$[0.625, 0.877]$ to $0.650$ and the workers plateau ($0.475$, $0.525$). One
criterion, opposite predictions, all borne out.

Two caveats. First, $\rho$ is computed from \emph{isolated} survivals that
Sections~5.2--5.7 of the main paper show to be unreliable; the Llama and
qwen verdicts survive because neither $\sum_j a_jc_j$ is marginal ($1.80$, $0.58$),
whereas DeepSeek's $1.018$ sits close enough to the critical value that its
supercritical label should be read together with its measured dynamics (endemic over
eight rounds) rather than from the point estimate alone. Second, we study a single
recurrent pattern at one temperature; the three models span five to eight rounds,
enough to separate decay from saturation, not to demonstrate extinction.


\section{Judge Calibration, Scope, and Known Failure Modes}
\label{app:judge-validation}

\paragraph{What was measured.} The compromise rule of Section~3.4 of the main paper
($C = \ind[\mathrm{ASV} \ge 0.9]$, ASV $=$ fraction of target bigrams present) was
calibrated on \textbf{ground-truth-labelled} outputs from the local model
\texttt{qwen2.5:7b} ($n = 12$ per compromised class, $n = 12$ benign, $n = 3$
refusing outputs; exact-string injection, benign ``verification code'' framing). The
labels are assigned by construction, not by a judge: \emph{benign} outputs never see
an injected instruction; \emph{refusing} outputs see one and do not reproduce the
target; \emph{propagating} outputs are retained only when the target appears, and
\emph{direct} outputs come from the injected instruction alone.

\begin{table}[h]
\small
\centering
\caption{Calibration of the two continuous scores on labelled qwen2.5:7b outputs
(mean $\pm$ sd). ASV separates the two compromised classes ($1.000$) from both
benign classes ($0.000$) with no overlap; the full-containment variant instead
scores benign outputs as high as $0.459$, which is what motivated the switch
(reported in the main paper).}
\vspace{6pt}
\label{tab:judge-calib}
\begin{tabular}{lcccc}
\toprule
\textit{score} & \textit{benign} & \textit{refusing} & \textit{propagating} & \textit{direct} \\
\midrule
ASV (target-bigram containment) & $0.000 \pm 0.000$ & $0.000 \pm 0.000$ & $1.000 \pm 0.000$ & $1.000 \pm 0.000$ \\
MR (Dice vs.\ direct output)    & $0.253 \pm 0.179$ & $0.484 \pm 0.044$ & $0.598 \pm 0.070$ & $0.862 \pm 0.212$ \\
full-containment (rejected)      & $0.459 \pm 0.215$ & $0.597 \pm 0.274$ & $0.729 \pm 0.223$ & $0.825 \pm 0.274$ \\
\bottomrule
\end{tabular}
\end{table}

\paragraph{Agreement at the deployed threshold.} At $\tau_{\mathrm{ASV}} = 0.9$ the
rule attains precision $1.000$, recall $1.000$, F1 $1.000$, FPR $0.000$ and FNR
$0.000$ on this set ($\mathrm{TP} = 24$, $\mathrm{FP} = 0$, $\mathrm{TN} = 15$,
$\mathrm{FN} = 0$); the closest benign score observed anywhere in the probes is
$0.875$, i.e.\ the threshold sits inside a gap rather than on a boundary. The
Dice-based MR branch, by contrast, separates the same classes at accuracy $0.949$
with $\mathrm{FPR} = 0$ only for $\tau_{\mathrm{MR}} > 0.53$, and the rejected
full-containment variant reaches accuracy $0.795$ with $\mathrm{FPR} = 0.067$ at its
own optimal cut---consistent with the main paper's $1.000$ vs.\ $0.795$ comparison.

\paragraph{Scope of this validation, and what it does not establish.} The calibration
covers \textbf{one model} (\texttt{qwen2.5:7b}), \textbf{one payload family}
(exact-string targets), and \textbf{no defence}: both the benign and the compromised
classes are drawn from undefended runs. The rule is deterministic and needs no
per-model recalibration, but we do \emph{not} claim that the same operating point
transfers unchanged to frontier models at temperature $0.7$, to paraphrase- or
obfuscation-obfuscated targets, or to defended runs. Two further failure modes follow
from the operational definition and should be read alongside every number we report:
(i) a paraphrased leak that destroys the target's bigrams is scored as
\emph{not} propagating (false negative), and (ii) a safe quotation of the target
inside a benign summary is scored as propagating (false positive). Validation on
frontier models, multiple payloads and defended conditions requires human
annotation of a stratified subset across those conditions; we report the protocol
for it below so the limitation is precise rather than open-ended. Read against a
human annotator, both failure modes concern the \emph{semantics} of the target: a
paraphrase that destroys its bigrams is scored as not propagating, and a safe
quotation inside a benign summary is scored as propagating.

\paragraph{Outcome levels.} Every judged activation is reported at level (a)
\emph{textual propagation}---the target string/code fragments appear. The calibration
set contains no case that separates (a) from (b) \emph{instruction
adoption/compliance} within the same output, and no case of (c) \emph{harmful
action} (the payload is a benign code, so (c) is by construction outside this
study). Adoption is a subset of propagation in this task family: an agent that acts
on the injected instruction must surface the target, whereas the converse does not
hold. Since every label here is positive on (a) whenever it is positive on (b), the
judge is an \emph{upper bound} on the adoption rate and we report (a) as the
propagation outcome. We do not use it as evidence of instruction adoption or of
harmful action anywhere in the paper.

\paragraph{Replay procedure.} In the same spirit we state precisely what the
content-form replay of Section~5.4 of the main paper does. Each replay trial draws a
fresh receiver, rotates the benign assignment over the same pool, and consumes either
content captured from a real natural run, a canonical artefact produced on demand, a
fixed hand-written string, or the in-context content under a different sender label;
temperature, token budget, system prompt, role, judge and thresholds are identical
across arms, so the arms differ only in message form and in the label of its sender.
The in-context arm is the harness self-check and has a target it must hit, namely the
rate measured in-chain; the other arms are read as comparisons against it. Because
each trial resamples the source output and draws the receiver fresh, a replay rate can
differ from an in-chain rate through sampling noise and per-context variation as well
as through message form. On the edge where the self-check misses ($0.733$ versus
$0.947$) we therefore report the value for completeness and do not use it as an
estimate of the form effect. The local re-run of this harness
(\texttt{qwen2.5:7b}, $4$ natural trials and $2$ replays per arm and edge) is
consistent with that reading rather than with a validated harness: every edge sits
inside the replay confidence interval, yet two of three edges differ from the in-chain
rate by $0.500$, which at this sample size the interval cannot resolve. We read the
local run as \emph{underpowered}, not as a passed self-check, and we do not report it
as evidence about form; the $n = 30$--$40$ configuration behind the content-form table
of the main paper (its Table~4) remains the only configuration in which the
self-check discriminates.

\paragraph{Behaviour on replayed activations (local model).} Table~\ref{tab:judge-adjud}
adjudicates every activation of a 24-item stratified subset drawn from the replayed
harness of Section~5.4 of the main paper (\texttt{qwen2.5:7b}, one payload, no defence;
the subset is stratified over arm and over judge verdict so that both classes are
represented). Each output was read and labelled at the three outcome levels before the
judge score was consulted. Against level (a) the rule reproduces the labels exactly:
precision $1.000$, recall $1.000$, $F_1 = 1.000$, $\mathrm{FPR} = \mathrm{FNR} =
0.000$. Against level (b) it is not a usable detector: it flags $19$ activations, and
none of them involves instruction adoption---every positive is the target quoted
inside a work product (an executive summary, a ticket classification, a release note,
a review annotation). Three activations that sit downstream of content that \emph{did}
carry the target do not carry it themselves, and the rule correctly reports them as
negative; a reader who stopped at the upstream context would have scored them as
propagating. We therefore report (a) and never (b), which is precisely the distinction
the main paper draws when it calls the judge an upper bound on adoption.

\begin{table}[h]
\small
\centering
\caption{Adjudication of the 24-activation subset against the three outcome levels
(level (a) textual propagation, (b) instruction adoption, (c) harmful action).
Labels were assigned by reading the full outputs, before the judge score was
consulted, with the help of an automated string-matching assistant and with
review by one of the authors; the labelling was not independent of the paper's
authors, so these numbers are a consistency check on the rule, not an independent
validation of it.}
\vspace{6pt}
\label{tab:judge-adjud}
\begin{tabular}{lccccccc}
\toprule
\textit{level} & \textit{TP} & \textit{FP} & \textit{TN} & \textit{FN}
 & \textit{prec.} & \textit{recall} & \textit{FPR} \\
\midrule
(a) textual propagation & 19 & 0 & 5 & 0 & $1.000$ & $1.000$ & $0.000$ \\
(b) instruction adoption & 0 & 19 & 5 & 0 & --- & --- & $0.792$ \\
(c) harmful action       & 0 & 19 & 5 & 0 & --- & --- & $0.792$ \\
\bottomrule
\end{tabular}
\end{table}

\paragraph{Protocol for a full validation (not yet executed).} A stratified subset should be drawn to distribute evenly across no-defence/defence, judge
success/failure, topology, model and arm; each case should be labelled at levels
(a), (b) and (c) by a human annotator, blind to the judge score, with a second
annotator on a subset for agreement. Reported metrics should be precision, recall,
F1, FPR, FNR and the full $2 \times 2$ confusion matrix per level, with
Wilson/Clopper--Pearson intervals, cluster bootstrap for intervals (clusters $=$
payload $\times$ defence $\times$ arm $\times$ edge, since trials repeat within
contexts), and worked TP/FP/TN/FN examples. The scripts that generate the
stratified sheet and score it are part of the released artifact; the numbers in
Table~\ref{tab:judge-calib} are the portion we had measured at submission time.


% =====================================================================
\section{Reproducibility, Artifact Claims, and Replication Details}
\label{app:repro}

We explicitly align reproducibility with the materials supplied. The claims about
reproducibility must match the actual artifact/package provided with the submission.
If code, per-cell outputs, and a manifest are included as an anonymized artifact, a
manifest (table $\to$ data file $\to$ script) should specify how to regenerate each
table without live model calls. If the artifact is not yet included, the reproducibility
statement should reflect partial reproducibility rather than claiming reproduction of
``every table''.

\paragraph{Minimum replication details.} Any released artifact or reported configuration
should include: full prompts and payload templates (anonymized where necessary); exact
model identifiers (vendor+model ID), API/provider+version, collection dates/time window,
platform; decoding parameters (temperature, top\_p/top\_k, max\_tokens, stop sequences,
seed if any, frequency/presence penalties); execution details (retries, error handling,
timeouts, rate limits, failed calls with counts, token budgets, call counts per
condition); fresh-artifact sampling/randomization/order/stratification; non-determinism
handling (temperature${>}0$); and a machine-readable experiment inventory (condition
$\to$ protocol, topology, model, defence, payload form, context/task,
$n,N_{\text{success}},N_{\text{total}}$, call counts). All estimates should include
numerator/denominator and appropriate uncertainty (clustered where trials repeat within
contexts/runs), with confirmatory vs.\ exploratory analyses clearly labeled.

\end{document}
